\documentclass[11pt]{article}

\usepackage[T1]{fontenc}
\usepackage[utf8]{inputenc}
\usepackage{lmodern}
\usepackage{amsmath,amssymb,amsfonts,amsthm}
\usepackage[margin=1in]{geometry}
\usepackage{microtype}
\usepackage{xcolor}
\usepackage[colorlinks=true,linkcolor=teal,urlcolor=teal]{hyperref}

\newcommand{\T}{\mathcal T_1}
\newcommand{\F}{\mathcal F_{\mathrm{orb}}}
\newcommand{\Tr}{\operatorname{Tr}}
\newcommand{\dd}{\,\mathrm d}
\newcommand{\B}{\mathrm B}

\newtheorem*{lemmaC}{Lemma C.1 (rephrased)}

\title{A Short Guide to Lemma C.1\\
\large Weighted root-fidelity bounds}
\author{}
\date{}

\begin{document}
\maketitle

\begin{abstract}
Lemma C.1 converts root fidelity into a product of weighted moments.  Its
first part is an operator inequality.  Its second part combines that
inequality with classical Cauchy--Schwarz for a classical--quantum state.
The weights used later are chosen so that one moment is compact and can pass
to a weak limit, while the remaining factors are exactly the classical and
orbital Gaussian affinities.  This note explains the regularized inverse,
the two Cauchy--Schwarz steps, and the role of the lemma in Proposition C.5.
\end{abstract}

\section{Why a weighted bound is useful}

For positive trace-class operators $R$ and $T$, their root fidelity is
\[
  \B(R,T)=\|R^{1/2}T^{1/2}\|_1.
\]
Appendix C must control this nonlinear quantity along sequences that converge
only after testing against compact observables.  Fidelity itself is not one
of those linear tests.  Lemma C.1 replaces it by a moment
\[
  R\longmapsto\Tr(RW)
\]
for a bounded positive compact operator $W$.  This moment does pass to the
limit supplied by Lemma C.2.

The inverse weight is paid on the fixed target:
\[
  \B(R,T)^2
  \leq \Tr(RW)\Tr(TW^{-1}).
\]
Thus $W$ should make the first factor continuous in the available topology
and the second factor explicitly computable.  The witness $W_\gamma$ used in
Appendix C has exactly these two properties.

\section{The regularized inverse}

The weight $W$ is bounded, positive, and injective, but it need not be
bounded below.  Its inverse may therefore be unbounded.  The expression
$\Tr(TW^{-1})$ is defined through bounded regularizations:
\begin{equation}\label{eq:inverse-definition}
  \Tr(TW^{-1})
  :=\lim_{\varepsilon\downarrow0}
  \Tr\bigl(T(W+\varepsilon I)^{-1}\bigr).
\end{equation}
As $\varepsilon$ decreases, the positive operators
$(W+\varepsilon I)^{-1}$ increase in the spectral order.  Hence the limit
exists in $[0,\infty]$.  The operator inequality is useful when this limit is
finite.

Injectivity rules out a genuine zero eigenspace.  It does not make $W^{-1}$
bounded, and that is why the definition in \eqref{eq:inverse-definition} is
needed.  For the later thermal witness, the eigenvalues of $W_\gamma$ tend to
zero, so this distinction is essential.

\section{Statement of Lemma C.1}

\begin{lemmaC}
Let $R,T$ be positive trace-class operators, and let $W$ be bounded,
positive, and injective.  If \eqref{eq:inverse-definition} is finite, then
\begin{equation}\label{eq:quantum-bound}
  \B(R,T)^2\leq\Tr(RW)\Tr(TW^{-1}).
\end{equation}

Let $\mathsf E$ be a Euclidean space, let
$R\in L^1(\mathsf E;\T(\F))$ be positive, let $g$ be a probability density,
and let $\chi>0$ be measurable.  Whenever the right side is finite,
\begin{equation}\label{eq:hybrid-bound}
  \B(R,gT)^2
  \leq
  \left(\int_{\mathsf E}\chi(y)\Tr[R(y)W]\dd y\right)
  \left(\int_{\mathsf E}\frac{g(y)}{\chi(y)}\dd y\right)
  \Tr(TW^{-1}).
\end{equation}
\end{lemmaC}

The notation $gT$ means the classical--quantum density
$y\mapsto g(y)T$.  When $\dim\mathsf E=0$, the classical integrals reduce to
evaluation at the unique point, and \eqref{eq:hybrid-bound} becomes
\eqref{eq:quantum-bound}.

\section{Proof walkthrough}

\subsection*{Step 1: insert a bounded inverse}

Set
\[
  W_\varepsilon=W+\varepsilon I.
\]
This operator is strictly positive and has a bounded inverse.  Insert
$W_\varepsilon^{1/2}W_\varepsilon^{-1/2}=I$ between the two square roots:
\[
  R^{1/2}T^{1/2}
  =R^{1/2}W_\varepsilon^{1/2}
   W_\varepsilon^{-1/2}T^{1/2}.
\]
Schatten Cauchy--Schwarz, namely
$\|AB\|_1\leq\|A\|_2\|B\|_2$, gives
\begin{align}
  \B(R,T)
  &\leq
  \|R^{1/2}W_\varepsilon^{1/2}\|_2
  \|W_\varepsilon^{-1/2}T^{1/2}\|_2,\notag\\
  \B(R,T)^2
  &\leq
  \Tr(RW_\varepsilon)
  \Tr(TW_\varepsilon^{-1}).
  \label{eq:regularized-bound}
\end{align}
The identities in the second line follow from the definition of the
Hilbert--Schmidt norm and cyclicity of the trace.

\subsection*{Step 2: remove the regularization}

The first factor in \eqref{eq:regularized-bound} satisfies
\[
  \Tr(RW_\varepsilon)=\Tr(RW)+\varepsilon\Tr R
  \longrightarrow\Tr(RW).
\]
The second factor converges by the definition of $\Tr(TW^{-1})$.  Letting
$\varepsilon\downarrow0$ proves \eqref{eq:quantum-bound}.

\subsection*{Step 3: apply the operator bound pointwise}

For classical--quantum positive densities, root fidelity decomposes over the
classical variable.  Since the target density is $g(y)T$,
\[
  \B(R,gT)
  =\int_{\mathsf E}\sqrt{g(y)}\,\B(R(y),T)\dd y.
\]
Apply \eqref{eq:quantum-bound} to $R(y)$ and $T$:
\[
  \B(R(y),T)
  \leq
  \sqrt{\Tr[R(y)W]}
  \sqrt{\Tr(TW^{-1})}.
\]
The target factor is independent of $y$, so
\begin{equation}\label{eq:before-classical-CS}
  \B(R,gT)
  \leq
  \sqrt{\Tr(TW^{-1})}
  \int_{\mathsf E}
  \sqrt{g(y)\Tr[R(y)W]}\dd y.
\end{equation}

\subsection*{Step 4: split the classical weight}

Insert $\chi(y)^{1/2}\chi(y)^{-1/2}$ into the integrand in
\eqref{eq:before-classical-CS}:
\[
  \sqrt{g(y)\Tr[R(y)W]}
  =\sqrt{\chi(y)\Tr[R(y)W]}
   \sqrt{\frac{g(y)}{\chi(y)}}.
\]
Ordinary Cauchy--Schwarz on $\mathsf E$ gives
\begin{align*}
 &\left(\int_{\mathsf E}
  \sqrt{g(y)\Tr[R(y)W]}\dd y\right)^2\\
 &\qquad\leq
 \left(\int_{\mathsf E}\chi(y)\Tr[R(y)W]\dd y\right)
 \left(\int_{\mathsf E}\frac{g(y)}{\chi(y)}\dd y\right).
\end{align*}
Combining this with \eqref{eq:before-classical-CS} proves
\eqref{eq:hybrid-bound}.

\section{The witnesses used in Proposition C.5}

For the oscillator component, Appendix C takes
\[
  W_\gamma
  =\sum_{\mathbf k}
  \sqrt{\frac{t_{\mathbf k}}{t_{\mathbf k}^{(\gamma)}}}
  |\mathbf k\rangle\!\langle\mathbf k|.
\]
It is compact and injective, and it satisfies
\[
  \Tr(\tau_pW_\gamma^{-1})
  =\Tr(\tau_p^{(\gamma)}W_\gamma)
  =B_{\mathrm{orb}}(\gamma,p).
\]
For the full classical--quantum case, the classical weight is
\[
  \chi(y)=\sqrt{\frac{f_1(y)}{f_\gamma(y)}},
  \qquad
  \int_{\mathsf E}\frac{f_1(y)}{\chi(y)}\dd y
  =B_{\mathrm{cl}}(\gamma,d).
\]
Thus Lemma C.1 turns fidelity into one channel-dependent compact moment and
two fixed affinity factors.  Proposition C.3 bounds the quantum moment,
while Lemma C.4 controls the trace available to it in the full case.  For
the orbital case, $\mathsf E=\{0\}$ and $\chi=1$, so the classical factor is
absent.

\section{Points worth keeping straight}

\begin{enumerate}
  \item The inverse of $W$ need not be bounded.  Only its regularized
  expectation in the fixed target must be finite.
  \item Compactness of $W$ is not required for Lemma C.1 itself.  It is
  required later to pass the moment through the compact-observable limit.
  \item Two different Cauchy--Schwarz inequalities are used: Schatten
  Cauchy--Schwarz for operators and ordinary Cauchy--Schwarz for the
  classical variable.
  \item The same statement covers the orbital case by taking the classical
  Euclidean space to have dimension zero.
\end{enumerate}

\end{document}
